%%=============================================================================
%% Conclusie
%%=============================================================================

\chapter{\IfLanguageName{dutch}{Conclusie}{Conclusion}}%
\label{ch:conclusie}

\section{Answer to the research question}%
\label{sec:conclusion-answer}

This thesis started by investigating to what extent Meshtastic, as a decentralized \gls{lora}-based mesh network, offers a reliable and readily deployable emergency communication solution when primary network infrastructure fails.
Based on the initial literature study and the \gls{poc}, the answer to this question would be yes, on certain conditions.

The \gls{poc} showed that a low-cost Meshtastic setup can provide a real, working link over real distances without relying on any existing infrastructure. A direct connection between two nodes reached a maximum of 18.5\,km, and adding a single relay in this setup improved the end-to-end distance (from Base-Node to mobile node) to 30\,km, which is larger than the literature study suggested.
The \gls{poc} also proved a relay could convert a dead zone connection into a working connection (session 1 compared to session 3). When comparing against the five initial requirements the literature study gave us in chapter~\ref{ch:stand-van-zaken}, Meshtastic can be seen as a realistic alternative.
It performs well on infrastructure independence, deployability and affordability, and for security, while it has limitations, it offers a reasonable level of protection if correctly configured.
These figures come from three sessions in a single region with one hardware type, so they should be read as an indication rather than a guaranteed performance benchmark.

For the problem domain, the identified gaps were solutions that don't depend on main infrastructure, which are quick to deploy and are affordable at scale. Both \gls{aprs} and \gls{dmr} fail because of the licensing requirements. Satellite messengers like Garmin inReach and Zoleo mainly fail because of the recurring costs.
Meshtastic is the only alternative that ticks every box. Regarding the environmental factors, the \gls{poc} showed us that the trees, buildings and other obstructions impede the signal way more than just the distance or the altitude.
This was clearly illustrated in section~\ref{sec:poc-session2} where the highest test locations (Hotond, Kwaremont) had 0\,\% \gls{pdr} because of the trees, while lower altitude locations but with a clear \gls{los} achieved 100\,\%.
The security question is a bit tricky on Meshtastic. The messages are effectively encrypted using \gls{aes}-256 once a private key is configured, but the default channel is public, so on a wrong config every message isn't really encrypted. Also once set up, the channels don't use authentication to make sure the message is from the right device.
If someone gets the key, they can impersonate a node and send messages on the channel. The header of the message is also in cleartext, so the metadata is not protected. The private communication between devices (end-to-end, no channel usage) is protected by encryption and, unlike the shared channels, also by authentication.

The range and reliability question from the solution domain has also been answered in the \gls{poc}.
The maximum direct connection distance achieved between two nodes was 18.5\,km, and with a single relay, the end-to-end distance was 30\,km. The effective range of the relay was estimated at 10 to 13\,km depending on terrain.
On the security question, just as mentioned before, confidentiality holds up, but authentication and metadata protection do not. This has also been discussed in section~\ref{sec:poc-security}.

Compared to the other alternatives, Meshtastic is the only one that combines infrastructure independence, no licence requirement, affordability at scale and low-barrier deployment. The current setup is easy to set up and understand, but needs some technical know-how to install (section~\ref{sec:poc-config}).

\section{Reflection}%
\label{sec:conclusion-reflection}

The effective total range of a clear \gls{los} connection was surprising going into the \gls{poc}. The assumption was that the higher the node, the better the connection. But the results showed that a lower location with a better \gls{los} outperforms most high altitude nodes surrounded by trees.
This is important for a real-world deployment: picking a relay location isn't just ``find the highest point'', but it's ``finding the point with the best \gls{los} in the direction you need coverage''.

Something that makes sense, but wasn't expected either was the discovery of third-party nodes during testing, some as far as Amiens, roughly 130\,km away.
This was as said unexpected, but it was convincing that the mesh topology is actually working as intended: nodes that weren't configured together still find a way to route traffic between them.

While Meshtastic is cheap and doesn't need a licence, which checks the deployability box on paper, the difficulties encountered during the setup in section~\ref{sec:poc-config} show that some patience, and a basic technical knowledge is necessary (CLI config if web client unreliable, etc.).

\section{Limitations}%
\label{sec:conclusion-limitations}

Some limitations were also discovered. The only tested hardware was the Heltec WiFi LoRa 32 V4; \gls{aprs}, \gls{dmr}, satellite messengers were only evaluated through literature study. Also other ESP32 options exist but weren't tested. They should work in the same way this \gls{poc} was done.
The comparison combines measured data on one side with documented data on the other. All measurements were also done in a single region (around Ronse and Tournai) and during a dry summer period. The results may look different in a more flat region, or during winter time, etc.
Latency, which was listed as a metric in the original proposal, was intentionally left out of the measured scope (see section~\ref{sec:method-measurement}). It was just informally noted, seeing traceroutes and messages arrived within a few seconds, which is good enough for this kind of scenarios.
Range and packet delivery turned out to be the decisive factors for answering the research question.
As planned, the security analysis stayed documentation-based combined with some setup observations. No real traffic has been intercepted or decrypted, both because that would be illegal and because it fell outside the scope of this thesis.

\section{Future research}%
\label{sec:conclusion-future}

The most obvious next step for continuing this research is to expand the \gls{poc} with more nodes, and more relays. This would allow to test even more hops and see how far the range could be extended using more nodes.
Each node probably adds some packet loss. It would be useful to know exactly how much and how far (in terms of hops and distance) this type of network can be pushed before it stops being useful.

A rough estimation can be made on creating a longer multi-hop path. The first hop of the current setup (base to relay) reached 18.5\,km with a clear \gls{los}. The second hop (relay to mobile) reached an effective 10--13\,km, which gave us a 30\,km end-to-end distance.
If a second relay would be added at another elevated, clear-\gls{los} site, a comparable third hop of roughly 10--15\,km could be estimated. This would extend the theoretical reach of the chain to approximately 45--50\,km over three hops.
This remains an estimation; every additional hop adds latency and more chance of packet loss. If a second relay would be introduced, the hop limit setting should also be moved from 3 to a higher count.
An effective test with a second relay should be done to confirm the effective range. That's a clear direction for further research.

Another next step could be testing in different environmental conditions. This was supposed to be done during the \gls{poc}, but the weather stayed largely dry and calm throughout the \gls{poc} (20 July 2026 – 20 August 2026),
so no systematic weather comparison was possible during the sessions. Only one day allowed testing in the rain (section~\ref{sec:poc-rain}).
This does not mean the setup is viable in all weather conditions, but it shows that rain does not necessarily break the connection, and this should be confirmed in future research.
The same measurements could be repeated but at different times of the day, and most importantly, in different weather conditions (rain, wind, temperature swings) to see whether the change in conditions has a real-world effect on the stability of the connection.

Going more into the security side of Meshtastic could be another direction to dig into, especially authentication and metadata protection (section~\ref{sec:sota-security}).
Lastly, as mentioned in section~\ref{sec:poc-session3}, using higher-gain or directional antennas (such as Yagi, parabolic, etc.) instead of a small omnidirectional one could largely improve the quality and the range of the network. The broadcast aspect would likely be lost, but in return the range and quality could be extended.
This wasn't tested during this thesis due to cost and time constraints, but it is a realistic next step for anyone trying to push the research even further.


\section{Comparison matrix}%
\label{sec:conclusion-comparison}
\begin{table}[H]
  \centering
  \begin{tabular}{lcccc}
    \toprule
    \textbf{Requirement} & \textbf{APRS} & \textbf{DMR} & \textbf{Satellite} & \textbf{Meshtastic} \\
    \midrule
    Infrastructure-independent   & Yes     & Partial & Yes    & Yes             \\
    Low-barrier deployment       & No      & No      & Yes    & Yes             \\
    Affordable at scale          & Partial & Partial & No     & Yes             \\
    No licence required          & No      & No      & Yes    & Yes             \\
    Range and reliability        & Good    & Good    & Global & Good            \\
    \midrule
    Data throughput              & Low     & High    & Low    & Very low        \\
    Message type                 & Text    & Voice + data & Text & Text + telemetry \\
    Guaranteed delivery          & No      & Yes     & Yes    & No (best-effort) \\
    Independent of node density  & Yes     & Yes     & Yes    & No              \\
    \bottomrule
  \end{tabular}
  \caption[Comparison of off-grid communication alternatives.]{\label{tab:conclusion-comparison}Comparison of the reviewed off-grid communication alternatives.}
\end{table}

The completed matrix and conclusion confirms that Meshtastic is not a universal winner.
It trades throughput, voice capability and guaranteed delivery for cost, licence-free use and easy deployment requirements.
For the specific scenario of this thesis, a low-cost text-based network that survives infrastructure failure, that trade-off is favourable, but it would be a poor choice where high bandwidth, voice, or guaranteed delivery are required.

\section{Recommendations for practitioners}%
\label{sec:conclusion-recommendations}

After a \gls{poc} like this, anyone considering Meshtastic for a real emergency deployment should take some practical points into account.
First, node placement should be planned around \gls{los}, not just around altitude.
Second, \textbf{always} replace the default encryption key with a random private key before using the network for anything sensitive, using the default key is basically the same as not encrypting any of the traffic at all.
Third, if main power isn't guaranteed, the nodes can easily run on a solar panel and battery (with the battery taking over at nighttime). In this \gls{poc} this was tested with a combination of a 6\,W, 5\,V/1.2\,A solar panel and a 3.7\,V, 2,000\,mAh battery. These could also be upgraded to a larger panel and battery if needed, but this setup was more than enough to run a node continuously through the night.
Both the solar powered nodes have been powered during a month without having to intervene.
Finally, the web interface not always being reliable during the setup, the mobile app or the command-line interface should be used for configuration instead.

With these points taken into account, Meshtastic stays a genuinely usable, low-cost option for emergency communication, well within the reach of volunteer organisations or small teams without a big budget or specialized licences.